\section{Experimentos}
En esta sección se detalla la infraestructura experimental, el procesamiento de las imágenes, la calibración de los hiperparámetros de entrenamiento y la configuración del estudio de ablación.

\subsection{Conjunto de Datos y Preprocesamiento}
El conjunto de datos de textura de neumáticos contiene dos divisiones principales para entrenamiento y validación independiente:
\begin{itemize}
    \item \textbf{training\_data}: Utilizado para el ajuste de los pesos de la red convolucional. Las muestras son equilibradas dinámicamente mediante el muestreo ponderado.
    \item \textbf{testing\_data}: Utilizado para la evaluación del desempeño general del modelo sin realizar aumentación de datos ni balanceo de lotes.
\end{itemize}

El procesamiento de las imágenes de entrada se estructuró a través del pipeline \texttt{torchvision.transforms}:
\begin{enumerate}
    \item \textbf{Redimensionamiento:} Todas las imágenes se reescalaron a una dimensión fija de $224 \times 224$ píxeles. Esto permite estandarizar las dimensiones espaciales para que sean compatibles con el extractor convolucional de la arquitectura ResNet-50.
    \item \textbf{Aumentación Estocástica:} Para mejorar la generalización textural y la robustez del modelo ante diferentes ángulos de inspección física, se aplicó un volteo horizontal aleatorio (\texttt{RandomHorizontalFlip}) en las imágenes de entrenamiento.
    \item \textbf{Normalización Z-score:} Se normalizaron los canales R, G y B utilizando los valores de media y desviación estándar precalculados sobre el conjunto global ImageNet:
    \[
    \mu_{RGB} = [0.485, 0.456, 0.406], \quad \sigma_{RGB} = [0.229, 0.224, 0.225]
    \]
\end{enumerate}

\subsection{Configuración de Hiperparámetros}
Los entrenamientos se realizaron en un entorno con soporte de GPU Nvidia CUDA. La configuración detallada de los hiperparámetros experimentales incluye:
\begin{itemize}
    \item \textbf{Optimizador:} Se utilizó el optimizador Adam por su capacidad de adaptar las tasas de aprendizaje para cada parámetro basándose en estimaciones de momentos de primer y segundo orden.
    \item \textbf{Tamaño del lote (Batch Size):} Establecido en 64 imágenes para equilibrar la estabilidad de la estimación del gradiente con las restricciones de memoria de GPU, optimizado para una arquitectura T4 GPU.
    \item \textbf{Épocas de Entrenamiento:} Tanto el modelo básico \texttt{CustomCNN} como las configuraciones de \texttt{ResNet-50} fueron entrenados durante 50 épocas para asegurar una convergencia profunda de los clasificadores.
    \item \textbf{Comparación de Tasas de Aprendizaje:} Se evaluó la estabilidad y velocidad de convergencia de la arquitectura \texttt{ResNet-50} sintonizando con dos valores de learning rate diferentes: $2 \times 10^{-4}$ y $1 \times 10^{-4}$ (escalados linealmente con el tamaño de lote) durante un intervalo rápido de 5 épocas.
\end{itemize}

\subsection{Configuración del Estudio de Ablación}
El estudio de ablación se diseñó para medir de forma aislada el impacto de dos factores clave sobre el rendimiento del clasificador:
\begin{enumerate}
    \item \textbf{Arquitectura del Extractor de Características:} Comparando una arquitectura convolucional simple diseñada desde cero (\texttt{CustomCNN}) contra una arquitectura profunda y compleja basada en aprendizaje por transferencia (\texttt{ResNet-50}).
    \item \textbf{Función de Pérdida (Loss Function):} Evaluando el impacto de mitigar el desbalance de clases comparando la pérdida clásica de entropía cruzada binaria (\texttt{BCEWithLogitsLoss}) frente a la pérdida adaptativa re-ponderada (\texttt{Focal Loss}).
\end{enumerate}
